\chapter[Un processus collaboratif]{\label{III-B} L'unification des thésaurus~: un processus collaboratif}


\lettrine{A}{près} avoir esquissé les fondements conceptuels et techniques de la structuration de \gls{thesaurus} au \mae, il convient d’entreprendre le chantier de leur unification~: ce travail, loin d'être purement technique, engage en réalité l’institution dans un processus où l’accord des métiers, la concertation des usages et la réflexion sur les finalités du langage utilisé sont tout aussi décisifs que le choix des outils.



Il s’agit donc ici de comprendre comment les thésaurus du musée, jusqu'ici dispersés en silos, peuvent être rassemblés sans perdre la richesse de leurs nuances~: comment impliquer tout d'abord les acteurs de changer des habitudes et de mettre en place une gouvernance qui soit l'affaire non plus d'un métier mais de l'ensemble du \ac{dsc}. Nous proposerons ensuite un aperçu de la méthode qui a été déployée au \mae~ pour répondre à ces besoins. Ce sont ces tensions, ces innovations, ces impasses et ces ouvertures que nous avons choisi d’examiner dans cette partie~: non pour livrer un mode d’emploi, mais pour mettre en lumière la dynamique humaine et intellectuelle qui sous-tend toute réforme documentaire. En effet, ces travaux préliminaires ont permis de mettre en lumière que, plutôt qu'un défi technique, l’unification des thésaurus au \mae~ est avant tout un travail de dialogue, d’acculturation et de partage pour construire une mémoire commune dans les divers métiers de l'institution.



\input{parties/mainmatter/III/B/III_B_1.tex}



\input{parties/mainmatter/III/B/III_B_2.tex}



\bigskip

\bigskip

\bigskip



\lettrine{L}{e} processus d’unification des thésaurus tel qu’il a commencé à être déployé au \mae~ révèle la profondeur des enjeux humains et organisationnels qui président à toute réforme documentaire~: c’est en effet la capacité de l’institution à mobiliser ses agents, à fédérer ses savoirs, à instaurer un dialogue entre les usages et les normes et à soutenir l’effort sur la durée qui semble déterminante. Ces constats ont invité à prolonger la réflexion du stage~: comment accompagner la transformation et soutenir les acteurs dans la maîtrise de la masse documentaire ? C’est à l’examen des apports potentiels des technologies émergentes — et en particulier de l’\acf{ia} — que sera consacrée la partie suivante, pour interroger leur pertinence, leurs limites et leur capacité à servir une gouvernance documentaire à la mesure du patrimoine confié au musée.